\chapter{Integrimi Monte Carlo}
\begin{qellime}
Studenti formulon integralin si pritje; ndërton vlerësuesin Monte Carlo dhe gabimin e tij; përdor kampionim të rëndësisë, variabla antitetike dhe variabla kontrolli; trajton integrale shumëdimensionale dhe pacaktueshmëri në modele fizike.
\end{qellime}

\section{Integrali si pritje}
Për domen $D$ me vëllim $V_D$ dhe mostra uniforme $X_i$,
\begin{equation}
 I=\int_Df(\vect x)\dd\vect x=V_D\E[f(X)],
\qquad
 \hat I_N=\frac{V_D}{N}\sum_{i=1}^Nf(X_i).
\end{equation}
Vlerësuesi është i paanshëm dhe
\begin{equation}
 \operatorname{SE}(\hat I_N)=V_D\frac{\sigma_f}{\sqrt N}.
\end{equation}
Avantazhi kryesor është se rendi $N^{-1/2}$ nuk përkeqësohet drejtpërdrejt me dimensionin. Disavantazhi është konvergjenca e ngadaltë në dimension të ulët krahasuar me kuadraturat me rend të lartë.

\section{Shembulli gjeometrik i $\pi$}
Në katrorin njësi, probabiliteti që $x^2+y^2\le1$ është $\pi/4$. Prandaj
\begin{equation}
 \hat\pi_N=4\frac{N_{\mathrm{brenda}}}{N}.
\end{equation}
Ky shembull është intuitiv, por jo metodë efikase për llogaritjen e $\pi$; shërben për të demonstruar treguesit, ligjin binomial dhe intervalin e besimit.

\begin{figure}[ht]
\centering\includegraphics[width=.95\textwidth]{18_monte_carlo_pi.pdf}
\caption{Vlerësimi gjeometrik i $\pi$ dhe konvergjenca e vlerësuesit.}
\end{figure}

\section{Gabimi statistik}
Varianca vlerësohet nga mostrat,
\begin{equation}
 s_f^2=\frac1{N-1}\sum_{i=1}^N(f_i-\bar f)^2,
\qquad \widehat{\operatorname{SE}}(\hat I_N)=V_D\frac{s_f}{\sqrt N}.
\end{equation}
Raportimi $\hat I\pm\mathrm{SE}$ pa treguar $N$, metodën e kampionimit dhe korrelacionin është i paplotë. Për integrandë me variancë të pafundme, CLT standard mund të mos zbatohet.

\begin{figure}[ht]
\centering\includegraphics[width=.74\textwidth]{19_gabimi_monte_carlo.pdf}
\caption{Shkallëzimi $N^{-1/2}$ në një integral katërpërmasor.}
\end{figure}

\section{Reduktimi i variancës}
\subsection{Kampionimi i rëndësisë}
Nëse $p$ është dendësia e synuar dhe $q$ dendësia e kampionimit,
\begin{equation}
 I=\E_q\left[f(X)\frac{p(X)}{q(X)}\right].
\end{equation}
Zgjedhja ideale vendos më shumë mostra ku kontributi është i madh. Peshat duhet të monitorohen për degenerim.

\subsection{Variablat antitetike}
Për $U\sim\mathcal U(0,1)$, çiftet $U$ dhe $1-U$ janë të korreluara negativisht për funksione monotone. Mesatarja e çiftit mund të ketë variancë më të vogël se dy mostra të pavarura.

\subsection{Variablat kontrolli}
Nëse $g$ ka pritje të njohur,
\begin{equation}
 \hat I_c=\bar f-c(\bar g-\E g).
\end{equation}
Koeficienti optimal është $c^*=\Cov(f,g)/\Var(g)$. Metoda përdor korrelacionin me një model të thjeshtë për të korrigjuar modelin e shtrenjtë.

\section{Integrale shumëfishta dhe rajone të ndërlikuara}
Një tregues $\mathbf1_D(\vect x)$ e kthen integrimin mbi rajon të parregullt në pritje mbi një kuti. Efikasiteti bie kur $D$ zë fraksion shumë të vogël. Atëherë përdoret parametrizim i rajonit, kampionim i rëndësisë ose ndarje stratifikore.

Quasi-Monte Carlo zëvendëson rastësinë me vargje me diskrepancë të ulët. Për integrandë të rregullt mund të tejkalojë $N^{-1/2}$, por vlerësimi i gabimit kërkon randomizim ose analiza të tjera.

\section{Përhapja Monte Carlo e pacaktueshmërisë}
Kur dalja $Y=g(\vect X)$ varet jolinearisht nga hyrje të pacaktuara, kampionohen hyrjet sipas shpërndarjes së tyre dhe llogaritet $Y$ për çdo realizim. Kjo ruan jo-Gaussianitetin, kufijtë dhe korrelacionet. Shpërndarja e daljes jep kuantile dhe probabilitete rreziku, jo vetëm mesatare.

\begin{figure}[ht]
\centering\includegraphics[width=.74\textwidth]{20_zbritja_mjetit_monte_carlo.pdf}
\caption{Shpërndarja e shpejtësisë së uljes kur lartësia, shpejtësia fillestare dhe shtytja kanë pacaktueshmëri.}
\end{figure}

\begin{shembull}[Zbritja e mjetit hapësinor]
Për çdo realizim kampionohen masa, shtytja, lartësia dhe shpejtësia fillestare. Integruesi dinamik jep shpejtësinë e kontaktit. Probabiliteti $P(v_{\mathrm{kontakt}}<v_{\mathrm{kufi}})$ është madhësi më informative se një trajektore me parametra mesatarë.
\end{shembull}

\begin{lidhje}
Laboratori krahason kuadraturën dhe Monte Carlo-n, verifikon $N^{-1/2}$, përdor një teknikë reduktimi variance dhe zbaton përhapjen e pacaktueshmërisë në zbritjen e mjetit.
\end{lidhje}

\section*{Pyetje kontrolli}
\begin{enumerate}
\item Pse Monte Carlo bëhet më tërheqës në dimension të lartë?
\item Si zgjidhet një variabël kontrolli e mirë?
\item Pse simulimi me parametrat mesatarë nuk jep domosdoshmërisht daljen mesatare?
\end{enumerate}
\section{Algoritmet Monte Carlo në Python}
\subsection{Integruesi bazë me gabim standard}
\begin{lstlisting}[language=Python,caption={Integrimi uniform në një kuti shumëpërmasore.}]
def integral_mc(f, kufijte, N, seed=2026):
    rng = np.random.default_rng(seed)
    kufijte = np.asarray(kufijte, float)
    a, b = kufijte[:, 0], kufijte[:, 1]
    x = rng.uniform(a, b, size=(N, len(a)))
    fx = np.asarray(f(x), float)
    vellimi = np.prod(b - a)
    vleresimi = vellimi * fx.mean()
    se = vellimi * fx.std(ddof=1) / np.sqrt(N)
    return vleresimi, se

f = lambda x: np.exp(x.sum(axis=1))
I, se = integral_mc(f, [(0, 1)]*4, N=200_000)
print(f"I = {I:.6f} +/- {se:.6f}")
\end{lstlisting}

\subsection{Kontrolli i ligjit $N^{-1/2}$}
Për secilin $N$ kryhen $R$ simulime të pavarura. Llogaritet RMSE kundrejt një reference dhe përshtatet pjerrësia në koordinata logaritmike.
\begin{lstlisting}[language=Python]
Ns = 2**np.arange(8, 18)
rmse = []
for N in Ns:
    v = [integral_mc(f, [(0, 1)]*4, N, seed=s)[0]
         for s in range(40)]
    rmse.append(np.sqrt(np.mean((np.array(v) - I_ref)**2)))
pjerrtesia = np.polyfit(np.log(Ns), np.log(rmse), 1)[0]
\end{lstlisting}
Pjerrësia pritet afër $-1/2$. Devijimi mund të vijë nga numër i pamjaftueshëm përsëritjesh ose nga variancë e madhe e integrandit.

\subsection{Variabli kontrollues}
\begin{lstlisting}[language=Python,caption={Reduktimi i variancës me një madhësi të korreluar.}]
def kontrolli(fx, gx, Eg):
    cov = np.cov(fx, gx, ddof=1)
    c = cov[0, 1] / cov[1, 1]
    korrigjuar = fx - c*(gx - Eg)
    return korrigjuar.mean(), korrigjuar.std(ddof=1)/np.sqrt(len(fx)), c
\end{lstlisting}
Raportohet faktori i reduktimit të variancës $\Var(f)/\Var(f_c)$. Nëse korrelacioni është i dobët, metoda mund të mos sjellë përfitim.

\subsection{Pacaktueshmëria e zbritjes së mjetit}
Algoritmi i plotë është i folezuar: jashtë kampionohen parametrat; brenda integrohet dinamika për çdo realizim. Rezultati është shpërndarja e shpejtësisë së kontaktit dhe probabiliteti i tejkalimit të kufirit.
\begin{lstlisting}[language=Python]
rng = np.random.default_rng(12)
parametra = rng.multivariate_normal(mu, Sigma, size=20_000)
v_kontakt = np.array([simulo_zbritjen(p) for p in parametra])
rreziku = np.mean(v_kontakt < v_kufi)
se_rreziku = np.sqrt(rreziku*(1-rreziku)/len(v_kontakt))
\end{lstlisting}

\begin{kujdes}
Në shqip dallohen \emph{vlerësuesi} (estimator), \emph{vlerësimi} (estimate), \emph{varianca} dhe \emph{gabimi standard}. Shprehja ``gabimi Monte Carlo'' duhet të shoqërohet me përkufizimin konkret.
\end{kujdes}

Literatura shqip për Monte Carlo është më e kufizuar; për terminologjinë probabilitare dhe të rizgjedhjes janë përdorur \cite{butka,tasho,hyka}, ndërsa baza algoritmike mbështetet te \cite{newman,robert,press}.

\subsection{Integrimi Monte Carlo si vlerësim i një pritjeje}
Le të jetë $X\sim p$ dhe
\begin{equation}
 I=\int h(x)p(x)\dd x=\E_p[h(X)].
\end{equation}
Vlerësuesi bazë është
\begin{equation}
 \widehat I_N=\frac1N\sum_{i=1}^{N}h(X_i).
\end{equation}
Për mostra të pavarura,
\begin{equation}
 \E[\widehat I_N]=I,
 \qquad
 \Var(\widehat I_N)=\frac{\sigma_h^2}{N}.
\end{equation}
Gabimi standard vlerësohet me $s_h/\sqrt N$. Teorema qendrore kufitare jep
\begin{equation}
 \frac{\widehat I_N-I}{s_h/\sqrt N}\approx\mathcal N(0,1).
\end{equation}
Kështu për ta përgjysmuar gabimin duhen afërsisht katër herë më shumë mostra. Shkalla $N^{-1/2}$ nuk varet drejtpërdrejt nga dimensioni, por varianca e integrandit dhe vështirësia e kampionimit mund të rriten fort me dimensionin.

Për integralin $J=\int_D f(\vect x)\dd\vect x$ mbi një rajon me vëllim $V_D$, nëse $X$ është uniforme në $D$, atëherë $J=V_D\E[f(X)]$. Metoda gjeometrike hit-or-miss përdor një ndryshore treguese dhe zakonisht ka variancë më të madhe se mesatarja e drejtpërdrejtë e $f$.

\subsection{Variablat antitetike}
Nëse $U\sim\mathcal U(0,1)$, edhe $1-U$ është uniforme. Vlerësuesi
\begin{equation}
 \widehat I_{\mathrm{anti}}=\frac1N\sum_{i=1}^{N}
 \frac{h(U_i)+h(1-U_i)}{2}
\end{equation}
ka variancë
\begin{equation}
 \frac{1}{2N}\left[\Var(h(U))+\Cov(h(U),h(1-U))\right].
\end{equation}
Për $h$ monotone, kovarianca zakonisht është negative dhe varianca ulet. Metoda nuk është automatikisht e dobishme për funksione jo monotone; kovarianca duhet matur.

\subsection{Variabli kontrollues}
Le të jetë $Y$ madhësia e interesit dhe $C$ një madhësi e korreluar me pritje të njohur $\mu_C$. Vlerësuesi
\begin{equation}
 Y_\beta=Y-\beta(C-\mu_C)
\end{equation}
mbetet i paanshëm. Varianca është
\begin{equation}
 \Var(Y_\beta)=\Var(Y)+\beta^2\Var(C)-2\beta\Cov(Y,C).
\end{equation}
Derivimi ndaj $\beta$ jep optimumin
\begin{equation}
 \beta_*=\frac{\Cov(Y,C)}{\Var(C)}.
\end{equation}
Në praktikë $\beta_*$ vlerësohet nga një mostër pilot ose nga e njëjta mostër, duke raportuar qartë procedurën.

\subsection{Kampionimi i rëndësisë për ngjarje të rralla}
Kur integrali dominohet nga një rajon që $p$ e viziton rrallë, kampionimi i drejtpërdrejtë prodhon shumicën e kontributeve zero dhe disa pesha shumë të mëdha. Me $X\sim q$,
\begin{equation}
 I=\E_q\left[h(X)\frac{p(X)}{q(X)}\right].
\end{equation}
Propozimi ideal formal është $q_*(x)\propto|h(x)|p(x)$, por varet nga integrali i panjohur. Ai shërben si parim projektimi: $q$ duhet të imitojë madhësinë e integrandit, jo vetëm dendësinë bazë.

\subsection{Përhapja Monte Carlo në zbritjen e mjetit hapësinor}
Në një model të thjeshtuar vertikal,
\begin{align}
 \dot h&=v,\\
 m\dot v&=-mg(h)+\tfrac12\rho(h)C_DA,v|v|+T(t),
\end{align}
ku $h$ është lartësia dhe $v<0$ gjatë zbritjes. Parametrat $C_D$, masa, dendësia atmosferike dhe shtytja kanë pacaktueshmëri. Çdo realizim përdor një mostër të parametrave, integron ODE-në dhe regjistron shpejtësinë në kontakt. Probabiliteti i dështimit vlerësohet nga treguesi $I_r=1$ kur $|v_{\mathrm{kontakt}}|$ tejkalon kufirin.

Një analizë bindëse raporton intervalin e besimit të probabilitetit, jo vetëm numrin e dështimeve. Për ngjarje shumë të rralla, simulimi i drejtpërdrejtë është joefikas; nevojitet kampionim i rëndësisë ose ndarje e ngjarjeve, me korrigjimin përkatës të peshave.

\begin{lstlisting}[language=Python,caption={Vlerësues Monte Carlo me gabim standard.}]
import numpy as np

def monte_carlo(h, gjenero, N, fare=2026):
    rng = np.random.default_rng(fare)
    x = gjenero(rng, N)
    y = np.asarray(h(x), float)
    vleresimi = y.mean()
    gabimi_standard = y.std(ddof=1) / np.sqrt(N)
    return vleresimi, gabimi_standard
\end{lstlisting}

\subsection{Përmbledhje dhe ushtrime}
Monte Carlo prodhon një vlerësim dhe një procedurë të matjes së gabimit. Numri i shifrave duhet të jetë në përputhje me gabimin standard.
\begin{enumerate}
\item Nxirrni variancën e metodës hit-or-miss për një sipërfaqe me probabilitet $p$.
\item Provoni formulën e $\beta_*$ dhe zvogëlimin maksimal të variancës në terma të korrelacionit.
\item Krahasoni kuadraturën deterministe dhe Monte Carlo kur dimensioni rritet.
\item Ndërtoni një analizë konvergjence me shumë fara dhe përshtatni pjerrësinë në grafik log--log.
\item Vlerësoni probabilitetin e uljes së pasigurt të mjetit dhe jepni interval Wilson.
\end{enumerate}

\subsection{Kampionimi i shtresëzuar}
Le të ndahet domeni në rajone pa mbivendosje $D_k$ me probabilitete $p_k$. Integrali shkruhet
\begin{equation}
 I=\sum_{k=1}^{K}p_k\mu_k,
 \qquad \mu_k=\E[h(X)\mid X\in D_k].
\end{equation}
Nëse në shtresën $k$ përdoren $n_k$ mostra të pavarura,
\begin{equation}
 \widehat I=\sum_kp_k\overline h_k,qquad
 \Var(\widehat I)=\sum_k\frac{p_k^2\sigma_k^2}{n_k}.
\end{equation}
Me kosto të njëjtë për mostër dhe $\sum_kn_k=N$, minimizimi jep ndarjen Neyman $n_k\propto p_k\sigma_k$. Pra më shumë mostra vendosen në shtresa me peshë ose variancë të madhe. Kur $\sigma_k$ nuk njihet, përdoret një fazë pilot.

Edhe ndarja proporcionale $n_k=Np_k$ zakonisht nuk është më e keqe se kampionimi i thjeshtë, sepse heq luhatjen e rastësishme të numrit të pikave në çdo shtresë. Në një integral njëpërmasor, ndarja e intervalit në qeliza dhe zgjedhja e një pike të rastit për qelizë kombinon mbulim të rregullt me paanshmëri.

\subsubsection{Mesatarizimi kushtor}
Nëse $Y$ varet nga $X$ dhe një burim tjetër rastësie, atëherë
\begin{equation}
 \E[Y]=\E[\E(Y\mid X)].
\end{equation}
Zëvendësimi i $Y$ me $g(X)=\E(Y\mid X)$ ruan pritjen dhe, nga ligji total i variancës,
\begin{equation}
 \Var(g(X))\le\Var(Y).
\end{equation}
Kjo teknikë quhet Rao--Blackwellizim ose mesatarizim kushtor. Në një model transporti, nëse një pjesë e dinamikës mund të integrohet analitikisht për një trajektore të dhënë, është më mirë të përdoret pritja kushtore sesa të simulohet një hedhje shtesë.

\subsubsection{Sekuencat me mospërputhje të ulët}
Kuazi-Monte Carlo zëvendëson pikat e pavarura me sekuenca që mbulojnë kubin në mënyrë më uniforme, si Halton ose Sobol. Për funksione me variacion të kufizuar, pabarazia Koksma--Hlawka lidh gabimin me mospërputhjen e pikave. Shkallëzimi teorik mund t'i afrohet $N^{-1}(\log N)^d$, më i mirë se $N^{-1/2}$ për dimension efektiv të moderuar.

Pikat nuk janë të rastit, prandaj gabimi standard i zakonshëm nuk zbatohet. Randomizimi i sekuencës, për shembull scrambling, ruan mbulimin dhe lejon disa përsëritje të pavarura për vlerësimin e gabimit. Në dimension shumë të lartë, renditja e ndryshoreve ka rëndësi: dimensionet e para të një sekuence zakonisht kanë cilësi më të mirë dhe duhet t'u caktohen drejtimet më të rëndësishme.

\subsubsection{Ngjarjet e rralla}
Për një probabilitet $p\ll1$, vlerësuesi binomial ka gabim relativ
\begin{equation}
 \frac{\sqrt{p(1-p)/N}}{p}\simeq\frac{1}{\sqrt{Np}}.
\end{equation}
Për të marrë gabim relativ $10\%$ duhen rreth $100/p$ realizime. Nëse $p=10^{-8}$, kjo kërkon rreth $10^{10}$ realizime dhe është e papërshtatshme. Kampionimi i rëndësisë zhvendos shpërndarjen drejt dështimit dhe korrigjon me peshën e gjasës.

Një propozim tepër agresiv mund të japë pesha me variancë të pafundme ose disa mostra dominuese. Diagnostika përfshin madhësinë efektive të peshave, maksimumin e peshës së normalizuar dhe përsëritjen me propozime alternative. Rezultati i ngjarjes së rrallë nuk duhet raportuar pa këto kontrolle.

\subsubsection{Buxheti i gabimit në Monte Carlo}
Gabimi total mund të përfshijë diskretizimin e trajektores, animin e modelit, gabimin e kampionimit dhe gabimin nga vlerësimi i peshave. Rritja e $N$ ul vetëm komponentin e kampionimit. Nëse ODE-ja brenda çdo realizimi zgjidhet me hap $h$, një strategji e balancuar zgjedh $h$ që gabimi numerik të jetë dukshëm më i vogël, por jo shumë më i vogël, se gabimi Monte Carlo.

Për të kontrolluar këtë ndarje kryhen dy eksperimente: në një $N$ të madh ndryshohet $h$ për të matur animin numerik; në një $h$ të vogël ndryshohet $N$ dhe fara për të matur variancën. Raporti final jep të dy komponentët dhe koston në vlerësime modeli ose kohë procesori.

